% OpenStax Calculus Volume 2 -- Section 7.4 Exercises: Area and Arc Length in Polar Coordinates
% Exercises reproduced from OpenStax, licensed under CC BY 4.0.
% Source: https://openstax.org/books/calculus-volume-3/pages/1-4-area-and-arc-length-in-polar-coordinates
\documentclass{exercises}
\title{OpenStax Calculus Volume 2}
\subtitle{Section 7.4 Exercises: Area and Arc Length in Polar Coordinates}
\sourceurl{https://openstax.org/books/calculus-volume-3/pages/1-4-area-and-arc-length-in-polar-coordinates}
\author{}
\date{}

\begin{document}
\maketitle


For the following exercises, determine a definite integral that represents the area.

\begin{enumerate}[start=1]
\item Region enclosed by $r=4$
\item Region enclosed by $r=3\sin \theta$
\item Region in the first quadrant within the cardioid $r=1+\sin \theta$
\item Region enclosed by one petal of $r=8\sin(2\theta)$
\item Region enclosed by one petal of $r=\cos(3\theta)$
\item Region below the polar axis and enclosed by $r=1-\sin \theta$
\item Region in the first quadrant enclosed by $r=2-\cos \theta$
\item Region enclosed by the inner loop of $r=2-3\sin \theta$
\item Region enclosed by the inner loop of $r=3-4\cos \theta$
\item Region enclosed by $r=1-2\cos \theta$ and outside the inner loop
\item Region common to $r=3\sin \theta$ and $r=2-\sin \theta$
\item Region common to $r=2$ and $r=4\cos \theta$
\item Region common to $r=3\cos \theta$ and $r=3\sin \theta$
\end{enumerate}

For the following exercises, find the area of the described region.

\begin{enumerate}[resume]
\item Enclosed by $r=6\sin \theta$
\item Above the polar axis enclosed by $r=2+\sin \theta$
\item Below the polar axis and enclosed by $r=2-\cos \theta$
\item Enclosed by one petal of $r=4\cos(3\theta)$
\item Enclosed by one petal of $r=3\cos(2\theta)$
\item Enclosed by $r=1+\sin \theta$
\item Enclosed by the inner loop of $r=3+6\cos \theta$
\item Enclosed by $r=2+4\cos \theta$ and outside the inner loop
\item Common interior of $r=4\sin(2\theta)$ and $r=2$
\item Common interior of $r=3-2\sin \theta$ and $r=-3+2\sin \theta$
\item Common interior of $r=6\sin \theta$ and $r=3$
\item Inside $r=1+\cos \theta$ and outside $r=\cos \theta$
\item Common interior of $r=2+2\cos \theta$ and $r=2\sin \theta$
\end{enumerate}

For the following exercises, find a definite integral that represents the arc length.

\begin{enumerate}[resume]
\item $r=4\cos \theta$ on the interval $0\le \theta \le \frac{\pi}{2}$
\item $r=1+\sin \theta$ on the interval $0\le \theta \le 2\pi$
\item $r=2\sec \theta$ on the interval $0\le \theta \le \frac{\pi}{3}$
\item $r=e^{\theta}$ on the interval $0\le \theta \le 1$
\end{enumerate}

For the following exercises, find the length of the curve over the given interval.

\begin{enumerate}[resume]
\item $r=6$ on the interval $0\le \theta \le \frac{\pi}{2}$
\item $r=e^{3\theta}$ on the interval $0\le \theta \le 2$
\item $r=6\cos \theta$ on the interval $0\le \theta \le \frac{\pi}{2}$
\item $r=8+8\cos \theta$ on the interval $0\le \theta \le \pi$
\item $r=1-\sin \theta$ on the interval $0\le \theta \le 2\pi$
\end{enumerate}

For the following exercises, use the integration capabilities of a calculator to approximate the length of the curve.

\begin{enumerate}[resume]
\item \techrequired $r=3\theta$ on the interval $0\le \theta \le \frac{\pi}{2}$
\item \techrequired $r=\frac{2}{\theta}$ on the interval $\pi \le \theta \le 2\pi$
\item \techrequired $r=\sin^{2}(\frac{\theta}{2})$ on the interval $0\le \theta \le \pi$
\item \techrequired $r=2\theta^{2}$ on the interval $0\le \theta \le \pi$
\item \techrequired $r=\sin(3\cos \theta)$ on the interval $0\le \theta \le \pi$
\end{enumerate}

For the following exercises, use the familiar formula from geometry to find the area of the region described and then confirm by using the definite integral.

\begin{enumerate}[resume]
\item $r=3\sin \theta$ on the interval $0\le \theta \le \pi$
\item $r=\sin \theta+\cos \theta$ on the interval $0\le \theta \le \pi$
\item $r=6\sin \theta+8\cos \theta$ on the interval $0\le \theta \le \pi$
\end{enumerate}

For the following exercises, use the familiar formula from geometry to find the length of the curve and then confirm using the definite integral.

\begin{enumerate}[resume]
\item $r=3\sin \theta$ on the interval $0\le \theta \le \pi$
\item $r=\sin \theta+\cos \theta$ on the interval $0\le \theta \le \pi$
\item $r=6\sin \theta+8\cos \theta$ on the interval $0\le \theta \le \pi$
\item Verify that if $y=r\sin \theta=f(\theta)\sin \theta$, then $\frac{dy}{d\theta}=f'(\theta)\sin \theta+f(\theta)\cos \theta$.
\end{enumerate}

For the following exercises, find the slope of a tangent line to a polar curve $r=f(\theta)$. Let $x=r\cos \theta=f(\theta)\cos \theta$ and $y=r\sin \theta=f(\theta)\sin \theta$, so the polar equation $r=f(\theta)$ is now written in parametric form.

\begin{enumerate}[resume]
\item Use the definition of the derivative $\frac{dy}{dx}=\frac{dy/d\theta}{dx/d\theta}$ and the product rule to derive the derivative of a polar equation.
\item $r=1-\sin \theta$; $(\frac{1}{2},\frac{\pi}{6})$
\item $r=4\cos \theta$; $(2,\frac{\pi}{3})$
\item $r=8\sin \theta$; $(4,\frac{5\pi}{6})$
\item $r=4+\sin \theta$; $(3,\frac{3\pi}{2})$
\item $r=6+3\cos \theta$; $(3,\pi)$
\item $r=4\cos(2\theta)$; tips of the leaves
\item $r=2\sin(3\theta)$; tips of the leaves
\item $r=2\theta$; $(\frac{\pi}{2},\frac{\pi}{4})$
\item Find the points on the interval $-\pi \le \theta \le \pi$ at which the cardioid $r=1-\cos \theta$ has a vertical or horizontal tangent line.
\item For the cardioid $r=1+\sin \theta$, find the slope of the tangent line when $\theta=\frac{\pi}{3}$.
\end{enumerate}

For the following exercises, find the slope of the tangent line to the given polar curve at the point given by the value of $\theta$.

\begin{enumerate}[resume]
\item $r=3\cos \theta$, $\theta=\frac{\pi}{3}$
\item $r=\theta$, $\theta=\frac{\pi}{2}$
\item $r=\ln \theta$, $\theta=e$
\item \techrequired Use technology: $r=2+4\cos \theta$ at $\theta=\frac{\pi}{6}$
\end{enumerate}

For the following exercises, find the points at which the following polar curves have a horizontal or vertical tangent line.

\begin{enumerate}[resume]
\item $r=4\cos \theta$
\item $r^{2}=4\cos(2\theta)$
\item $r=2\sin(2\theta)$
\item The cardioid $r=1+\sin \theta$
\item Show that the curve $r=\sin \theta \tan \theta$ (called a cissoid of Diocles) has the line $x=1$ as a vertical asymptote.
\end{enumerate}

\end{document}
